\subsection{RESERVES}

This rule explains the use of the "Reserve" and "Committed" sections of the Stacking Boxes.

\begin{enumerate}[label=(\alph*)]
  \item At the start of the first Round of Combat, both players divide their units in the involved Stacking Boxes, placing them as desired into either the "Reserve" or "Committed" sections, with the following restrictions:
  \begin{enumerate}[label=(\roman*)]
    \item All Wagon and Supply counters must be placed in the "Reserve" section.
    \item At least one-half of all Combat Factors must be placed in the "Committed" section.
  \end{enumerate}
  The Union Player turns his units upside-down while doing this to conceal his deployment from the Confederate Player. The Confederate Player, with his SCREEN, does not thave this problem.

  \item The Confederate Player now pulls out his "Comitted" units for the Union Player's inspection, and the Union places his "Committed" units back rightside-up for Confederate inspection.

  \item Combat is now handled normally, but only the "Committed" units are used. The "Reserve" units take no part in the Combat. The Phasing Player's "Reserve" units do lose Movement Allowance Factors for the Round of Combat.

  \item After the first segment of the first Round of Combat, both sides may shift units freely from their "Reserve" to their "Committed" area, or vice-versa, during their own segment of a Round only.

  \item If all "Committed" units are eliminated, this counts for Annhilation Points, even if there are surviving units in "Reserve". If all "Committed" units are forced to retreat, the "Reserve" units must retreat with them. However, if the "Committed" units suffer a "Disorder" result, this will not affect the "Reserve" units (unless, of course, the "Disorder" forces the "Committed" units to retreat).

  \item The non-lettered Stacking boxes provided on the Mapboard are for the use of units not represented by a lettered one. Simply note the location of the units on the Map, move them to one of the unlettered Stack Boxes, handle Combat, then return the units to the Map when done.
\end{enumerate}
